% Function descriptions for coin task analysis on non-padded sequences.
% Assumes math_commands.tex is loaded. Each paragraph is self-contained
% and suitable for inclusion in a NeurIPS/ICML/ICLR methods or appendix section.

\subsection{Coin Task}

We study a collection of coin-flipping tasks indexed by $k \in [K]$. For each
task and batch sample $b \in [B]$, let
\[
x^{(k,b)}_{1:t} = \bigl(x^{(k,b)}_1, \ldots, x^{(k,b)}_t\bigr)
\]
denote the length-$t$ prefix. From layer $l$ we extract the hidden state at
position $t$ and write it as $\vh^{(l,b)}_{k,t} \in \R^d$. When the layer index
is not essential we suppress it and write $\vh^{(b)}_{k,t}$.

\paragraph{Bayesian posterior over tasks.}
Each task $k$ is a categorical distribution over $V$ symbols with probability
vector $\vp_k = (p_k(0), \ldots, p_k(V{-}1)) \in \Delta^{V-1}$.
Because the coin task is i.i.d., the likelihood of the observed prefix
$x_{1:t}$ under task~$k$ factorises as
\[
  \P(x_{1:t} \mid Z = k)
  = \prod_{\tau=1}^{t} p_k(x_\tau)
  = \prod_{v=0}^{V-1} p_k(v)^{\,n_v^{(t)}},
\]
where $n_v^{(t)} = \sum_{\tau=1}^{t} \mathbf{1}[x_\tau = v]$ is the count of
symbol~$v$ in the first $t$ tokens.
The task pool contains $K_{\mathrm{tot}} = K_{\mathrm{maj}} + K_{\mathrm{min}}$
tasks with prior
\[
  \pi_k =
  \begin{cases}
    (1 - p_0) / K_{\mathrm{maj}} & k \in \{1, \ldots, K_{\mathrm{maj}}\}, \\
    p_0 / K_{\mathrm{min}}       & k \in \{K_{\mathrm{maj}}+1, \ldots, K_{\mathrm{tot}}\},
  \end{cases}
\]
where $p_0$ is the minor-task mixing probability.
Bayes' rule gives the posterior after $t$ tokens:
\begin{equation}
\label{eq:coin_posterior}
  \P(Z = k \mid x_{1:t})
  = \frac{\pi_k \prod_{v} p_k(v)^{\,n_v^{(t)}}}
         {\sum_{j=1}^{K_{\mathrm{tot}}} \pi_j \prod_{v} p_j(v)^{\,n_v^{(t)}}}.
\end{equation}
Taking logarithms, the unnormalised log-posterior is
\[
  \log \pi_k + \sum_{v=0}^{V-1} n_v^{(t)} \log p_k(v),
\]
so the posterior depends on the data only through the sufficient statistic
$\vn^{(t)} = (n_0^{(t)}, \ldots, n_{V-1}^{(t)})$.
In practice we compute~\eqref{eq:coin_posterior} incrementally via cumulative
counts and log-sum-exp normalisation for numerical stability.

\paragraph{Representation model.}
We adopt the additive decomposition
\begin{equation}
\label{eq:coin_rep_model}
\vh^{(b)}_{k,t}
= \vmu_t
+ \vtheta_{k,t}\!\bigl(x^{(k,b)}_t\bigr)
+ \vepsilon^{(b)}_{k,t}\!\bigl(x^{(k,b)}_{1:t}\bigr),
\end{equation}
where $\vmu_t$ is a position-dependent mean shared across tasks and tokens,
$\vtheta_{k,t}(\cdot)$ captures the joint effect of task identity and the
current token, and $\vepsilon^{(b)}_{k,t}$ is a residual that may depend on the
full history. Since the current token distribution depends on $k$ in the coin
task, task and token information are generally entangled in
$\vtheta_{k,t}(\cdot)$.

Our analysis focuses on two testable properties of~\eqref{eq:coin_rep_model}.

\paragraph{(P1) Vanishing residual variance.}
In the coin task, all information needed for next-token prediction is captured
by a sufficient statistic (equivalently, the Bayesian posterior), which becomes
increasingly stable as more tokens are observed. Accordingly, we expect the
residual to become batch-invariant once $(k, x_t)$ is fixed:
\begin{equation}
\label{eq:p1_limit}
\Var_b\!\Bigl[\vepsilon^{(b)}_{k,t}\bigl(x^{(k,b)}_{1:t}\bigr)\,\Big|\, x_t, k\Bigr]
\to 0
\qquad \text{as } t \to \infty .
\end{equation}

\emph{Empirical P1 variance.}
To estimate this conditional variability, for each $(t,k,v)$ we generate $B$
sequences from task $k$, set the current token to $x_t=v$, and extract the
corresponding hidden states $\{\vh^{(b)}_{k,t}\}_{b=1}^B$. Define the
conditional mean estimator
\[
\hat{\vmu}_{t,v,k} := \frac{1}{B}\sum_{b=1}^B \vh^{(b)}_{k,t},
\]
and the normalized P1 variance
\begin{equation}
\label{eq:p1var_def}
\widehat{\mathrm{P1Var}}(l,t)
:= \frac{
\frac{1}{|\gV_t|}\sum_{v\in\gV_t}\frac{1}{K}\sum_{k=1}^K
\frac{1}{B}\sum_{b=1}^B
\bigl\|\vh^{(l,b)}_{k,t}-\hat{\vmu}_{t,v,k}\bigr\|_2^2
}{
\frac{1}{|\gV_t|}\sum_{v\in\gV_t}\frac{1}{K}\sum_{k=1}^K
\bigl\|\hat{\vmu}_{t,v,k}\bigr\|_2^2
}.
\end{equation}
Low values of $\widehat{\mathrm{P1Var}}(l,t)$ indicate that $\vh^{(l,b)}_{k,t}$
is (nearly) determined by the pair $(k,x_t)$, consistent with~\eqref{eq:p1_limit}.

\Figref{fig:p1_variance} shows $\widehat{\mathrm{P1Var}}(l,t)$ versus position~$t$ for layers $l \in \{2,3,4,5\}$, computed on a model trained with $K_{\mathrm{maj}} = 3$ major tasks and $K_{\mathrm{min}} = 8$ minor tasks ($k = 3$, so $n_{\mathrm{minor}} = 2^3$), using $B = 32$ batch samples per (task, token, position) triple. At all layers, the normalised P1~variance decreases with position, confirming that the residual~$\vepsilon$ diminishes as the model observes more tokens. Later layers exhibit lower P1~variance throughout, indicating that deeper representations more rapidly converge to a function of $(x_t, k)$ alone.

\begin{figure}[ht]
  \centering
  \includegraphics[width=0.65\textwidth]{Figs/p1_variance_coin_3.png}
  \caption{Normalised P1 variance \eqref{eq:p1var_def} versus position $t$ for
  layers 2--5 when $K_{\text{min}} = 8$. The decrease toward zero supports (P1):
  residual variance diminishes as evidence accumulates, with deeper layers
  converging faster.}
  \label{fig:p1_variance}
\end{figure}

\paragraph{(P2) Linear task--token decomposition.}
Motivated by the belief-state hypothesis, we expect the hidden state to encode
both the current token and a posterior over tasks. Writing
\[
\alpha^{(b)}_{t,k'} := \P\!\left(Z=k' \mid x^{(k,b)}_{1:t}\right),
\]
a convenient sufficient condition is the existence of a (possibly learned) map
$g$ such that $g(\vtheta_{k,t}(x))$ recovers $(\alpha^{(b)}_{t,1:K},x)$.
If nonlinear interactions between task identity and the token are weak, this
suggests the approximation
\begin{equation}
\label{eq:linear_decomp}
\vtheta_{k,t}(x) \approx \vnu_{k,t} + \vphi_t(x),
\end{equation}
where $\vnu_{k,t}$ is task-specific (encoding the belief state) and $\vphi_t(x)$
depends only on the token. We refer to~\eqref{eq:linear_decomp} as a
\emph{linear task--token decomposition}. Its failure indicates genuinely
token-dependent task representations.

\emph{Empirical task-centered variance.}
Under~\eqref{eq:linear_decomp}, token effects cancel when centering across
tasks. Accordingly, we define the task-centered vector
\[
\vd^{(b)}_{k,t}
:= \vh^{(b)}_{k,t} - \frac{1}{K}\sum_{k'=1}^K \vh^{(b)}_{k',t},
\]
and its batch mean $\hat{\vd}_{k,t} := \frac{1}{B}\sum_{b=1}^B \vd^{(b)}_{k,t}$.
We measure the variability across batches via
\begin{equation}
\label{eq:taskvar_def}
\widehat{\mathrm{TaskVar}}(l,t)
:= \frac{
\frac{1}{K}\sum_{k=1}^K \frac{1}{B}\sum_{b=1}^B
\bigl\|\vd^{(l,b)}_{k,t}-\hat{\vd}_{k,t}\bigr\|_2^2
}{
\bigl\|\hat{\vmu}_t\bigr\|_2^2
},
\qquad
\hat{\vmu}_t := \frac{1}{KB}\sum_{k=1}^K\sum_{b=1}^B \vh^{(b)}_{k,t}.
\end{equation}
If~\eqref{eq:linear_decomp} holds and the residual in~\eqref{eq:coin_rep_model}
is small, then $\vd^{(b)}_{k,t}$ is approximately batch-invariant (equal to
$\vnu_{k,t}-\frac{1}{K}\sum_{k'}\vnu_{k',t}$), so
$\widehat{\mathrm{TaskVar}}(l,t)$ should approach zero as $t$ grows. Persistent
task-centered variance provides direct evidence of nonlinear task--token
interactions and/or substantial history-dependent residual structure.

\Figref{fig:task_variance} shows $\widehat{\mathrm{TaskVar}}(l,t)$ versus
position~$t$ for layers $l \in \{2,3,4,5\}$, computed on the same model as
\figref{fig:p1_variance} ($K_{\mathrm{maj}} = 3$, $K_{\mathrm{min}} = 8$,
$k = 3$), using $B = 64$ batch samples.
In the early layers, the task-centered variance is large and persists across
positions, indicating that the task--token interaction is not yet disentangled.
In later layers the variance decreases, consistent with the hidden
representation converging toward the linear decomposition~\eqref{eq:linear_decomp}.

\begin{figure}[ht]
  \centering
  \includegraphics[width=0.65\textwidth]{Figs/p2_variance_coin_3.png}
  \caption{Normalised task-centered variance~\eqref{eq:taskvar_def} versus
  position~$t$ for layers 2--5 when $K_{\text{min}} = 8$. Deeper layers exhibit
  lower variance, supporting the linear task--token
  decomposition~\eqref{eq:linear_decomp}.}
  \label{fig:task_variance}
\end{figure}



\paragraph{Historical perturbation analysis.}
To assess how much the model's next-token prediction relies on historical
hidden representations, we perform a layer-wise perturbation experiment.
For each task~$k$ in the pool (major and minor), we generate $B$ sequences
from that task's distribution.
At a target position~$t$ in layer~$l$, we leave the hidden state at
position~$t$ unchanged but randomly replace each historical hidden state
$\vh_l^{(j)}$ for $j \in \{0, \ldots, t{-}1\}$ with noise
$\vr \sim \gN(\vzero, \mI)$, independently with probability~$p$.
We then measure the change in next-token log-probability:
\[
  \Delta \log \P(x_{t+1})
    = \log \P_{\mathrm{pert}}(x_{t+1} \mid x_{0:t})
    - \log \P_{\mathrm{base}}(x_{t+1} \mid x_{0:t})\,,
\]
averaged over the $B$ batch samples and all $K$ tasks.
A value of $\Delta \log \P \approx 0$ indicates that the model does not rely on
historical hidden states at that layer for its prediction, suggesting that the
relevant information has already been compressed into the current position's
hidden state.
A strongly negative $\Delta \log \P$ indicates that historical representations
carry critical information at that layer.

\Figref{fig:hist_inject} shows $\Delta \log \P$ versus position~$t$ for
layers $l \in \{0,1,2,3,4\}$, with perturbation probability $p = 0.9$ and
$B = 64$ batch samples, evaluated at every 3rd position up to $t = 60$.
\figleft\ uses a model with $K_{\mathrm{min}} = 32$ ($k = 5$); \figright\ uses
$K_{\mathrm{min}} = 1024$ ($k = 10$).
In both cases, early layers (layer~0, 1) show the largest drop in
$\Delta \log \P$, indicating that they carry the most critical historical
information.
Later layers (layer~4, 5) are nearly unaffected, suggesting that by those
layers the model has compressed the relevant history into the current
position's representation.

\begin{figure}[ht]
  \centering
  \includegraphics[width=0.48\textwidth]{Figs/hist_inject_coin_k5.png}
  \hfill
  \includegraphics[width=0.48\textwidth]{Figs/hist_inject_coin_k10.png}
  \caption{Historical perturbation $\Delta \log \P(x_{t+1})$ versus
  position~$t$ for layers 0--4 with $p = 0.9$.
  \figleft\ $K_{\text{min}} = 32$ ($k = 5$).
  \figright\ $K_{\text{min}} = 1024$ ($k = 10$).
  Early layers carry more critical historical information; later layers are
  largely unaffected by perturbation.}
  \label{fig:hist_inject}
\end{figure}



\paragraph{Forward linear probe: predicting the posterior from hidden states.}
To test whether the hidden representation at a given layer encodes the Bayesian
posterior over tasks, we train a linear probe that maps hidden states to
posterior predictions.

\emph{Data collection.}
We draw $N = 2^{13}$ sequences from the training mixture (major and minor tasks
sampled uniformly) and run a single forward pass through the frozen model for
each batch of $B = 256$ sequences.
At layer~$l$, a forward hook on the attention block captures the residual-stream
output $\vh_{l,t} \in \R^D$ at every position $t \in \{5, 6, \ldots, 24\}$.
The ground-truth Bayesian posterior
$\vp^*_t = \P(Z = k \mid x_{1:t}) \in \Delta^{K_{\mathrm{tot}}-1}$
is computed from the known task distributions using cumulative token counts.
Because the coin task is i.i.d., the same posterior applies to all positions for
a given sequence; it is therefore replicated across the $P = 20$ extracted
positions.
This yields $N \times P$ training pairs $(\vh_{l,t},\, \vp^*_t)$.

\emph{Linear-softmax model.}
We fit
$\hat{\vp}(Z \mid \vh) = \softmax(\mW \vh + \vb)$
with $\mW \in \R^{K_{\mathrm{tot}} \times D}$, $\vb \in \R^{K_{\mathrm{tot}}}$,
minimising the KL divergence
\[
  \gL = \KL\bigl(\vp^* \,\|\, \hat{\vp}\bigr)
      = \sum_{k} p^*_k \log \frac{p^*_k}{\hat{p}_k}
\]
via Adam (learning rate $10^{-2}$, 200 epochs, mini-batch size 4096).
The data are split 80/20 into training and validation sets.
The validation KL loss measures how much posterior information is linearly
decodable from the hidden state at each layer.

\Figref{fig:forward_probe} shows the validation KL loss as a function of
$k = \log_2 K_{\mathrm{min}}$ for layers~0--3.
All layers track each other closely at small~$k$, where the posterior has few
components and is easy to decode. As the number of minor tasks grows, the loss
increases in~$k$. 
%with deeper layers consistently achieving
%lower loss. This confirms that later layers encode the task posterior more
%faithfully, even as the posterior becomes higher-dimensional.

\begin{figure}[ht]
  \centering
  \includegraphics[width=0.65\textwidth]{Figs/forward_probe_coin.png}
  \caption{Validation KL divergence of the forward linear probe
  (hidden state $\to$ task posterior) versus $k = \log_2 K_{\text{min}}$
  for layers~0--3.
  Each experiment uses $K_{\text{maj}} = 3$ major tasks and
  $K_{\text{min}} = 2^k$ minor tasks, with $B = 256$, $N = 2^{13}$ sequences,
  and positions $5$--$24$.
  Deeper layers achieve lower loss, indicating more linearly accessible
  posterior information.}
  \label{fig:forward_probe}
\end{figure}



\paragraph{Reverse linear probe: decomposing hidden states into interpretable
components.}
While the forward probe tests \emph{whether} posterior information is present,
the reverse probe quantifies \emph{how much} of the hidden-state variance each
information source explains.

\emph{Data collection.}
We draw $N = 2^{13}$ sequences in \emph{major-task mode} (sampling only from the
$K_{\mathrm{maj}} = 3$ major tasks) and extract, for each sequence at each
position $t \in \{40, 41, \ldots, 79\}$:
(i)~the hidden vector $\vh_{l,t} \in \R^D$ via a forward hook on the attention
block of layer~$l$,
(ii)~the model's output logits $\bm{\ell}_t \in \R^V$,
(iii)~the one-hot encoding of the current token
$\mathrm{onehot}(x_t) \in \{0,1\}^V$,
and (iv)~the ground-truth Bayesian posterior
$\vp^*_t = \P(Z = k \mid x_{1:t}) \in \Delta^{K_{\mathrm{maj}}-1}$
computed with \texttt{include\_minor=False} so that only the 3 major-task
components appear.
Because the coin task is i.i.d., the same posterior applies to all extracted
positions within a sequence; it is replicated across the $P = 40$ positions.
This yields $N \times P$ data pairs that are split 80/20 into training and
validation sets.

\emph{Closed-form linear regression.}
For each predictor $\vz \in \{\vp^*,\, \bm{\ell},\, \mathrm{onehot}(x_t)\}$,
we fit
$\hat{\vh} = \mW \vz + \vb$
via the pseudoinverse $[\mW; \vb] = (\tilde{\mZ}^\top \tilde{\mZ})^{-1}
\tilde{\mZ}^\top \mH$,
where $\tilde{\mZ}$ augments $\mZ$ with a column of ones, and report the
coefficient of determination
\[
  R^2 = 1 - \frac{\sum_i \|\vh_i - \hat{\vh}_i\|^2}
                  {\sum_i \|\vh_i - \bar{\vh}\|^2}
\]
on the held-out validation set.
A combined predictor that concatenates all three feature vectors is also fitted.
When the individual~$R^2$ values are approximately additive
($R^2_{\mathrm{combined}} \approx R^2_{\mathrm{post}} + R^2_{\mathrm{logit}}
  + R^2_{\mathrm{token}}$),
the three information components occupy approximately orthogonal subspaces
of the hidden representation.

\Figref{fig:reverse_probe} shows the validation~$R^2$ across layers~0--5 for
three experiments with increasing numbers of minor tasks.
When no minor tasks are present ($k = -1$, panel~a), the token one-hot
dominates in early layers while the posterior~$R^2$ grows steadily, reaching
$0.92$ at layer~5.
Logits contribute negligibly, consistent with the model not needing to
distinguish tasks at the output level.
With a small minor-task pool ($k = 2$, $K_{\mathrm{min}} = 4$, panel~b),
the token component decreases in later layers while the posterior component
rises more sharply, and the logits baseline remains near zero.
With a larger pool ($k = 8$, $K_{\mathrm{min}} = 256$, panel~c), the logits
baseline becomes substantial in deeper layers, nearly matching the
posterior~$R^2$ at layer~5 ($0.89$ vs.\ $0.99$).
This is expected: the Bayes-optimal next-token distribution is a linear
function of the posterior,
$\P(x_{t+1} \mid x_{1:t}) = \sum_k \P(Z{=}k \mid x_{1:t})\, p_k(x_{t+1})$,
so when the task pool is large and its distributions are diverse, this linear
map is nearly injective and the logits become an almost lossless encoding of
the posterior.
In the final layer the representation is being prepared for the output
projection, so the logits naturally capture the same posterior information.
Across all settings, the posterior~$R^2$ increases monotonically with depth,
confirming that later layers devote more representational capacity to encoding
the Bayesian posterior.

Panel~(d) of \figref{fig:reverse_probe} repeats the analysis using the
\emph{training mixture} (\texttt{sample\_mode="train"}) instead of major-task
mode, for $k = 2$ ($K_{\mathrm{min}} = 4$, $K_{\mathrm{tot}} = 7$).
The posterior now has all $K_{\mathrm{tot}}$ components and the data include
both major and minor task sequences.
The overall pattern is similar---posterior~$R^2$ grows with depth, token~$R^2$
declines---but two differences stand out.
First, the logits baseline rises sharply from layer~2 onward, reaching
$R^2 = 0.89$ at layer~5, nearly matching the posterior ($R^2 = 0.95$).
This contrasts with the major-only setting (panel~b), where the logits
baseline was essentially zero.
The explanation is that the training mixture contains minor tasks with distinct
categorical distributions; predicting the next token accurately requires the
model to encode which task is active, so the output logits carry substantial
posterior information even at moderate~$k$.
Second, the token~$R^2$ remains elevated through layer~1 ($0.52$) before
dropping, reflecting the fact that in the training mixture the token
distribution varies across tasks, so early layers retain more token-level
signal to help disambiguate tasks.

\begin{figure}[ht]
  \centering
  \begin{subfigure}[t]{0.48\textwidth}
    \centering
    \includegraphics[width=\textwidth]{Figs/reverse_probe_coin_k-1.png}
    \caption{Major mode, $k = -1$ (no minor tasks)}
  \end{subfigure}
  \hfill
  \begin{subfigure}[t]{0.48\textwidth}
    \centering
    \includegraphics[width=\textwidth]{Figs/reverse_probe_coin_k2.png}
    \caption{Major mode, $k = 2$ ($K_{\text{min}} = 4$)}
  \end{subfigure}

  \vspace{0.5em}
  \begin{subfigure}[t]{0.48\textwidth}
    \centering
    \includegraphics[width=\textwidth]{Figs/reverse_probe_coin_k8.png}
    \caption{Major mode, $k = 8$ ($K_{\text{min}} = 256$)}
  \end{subfigure}
  \hfill
  \begin{subfigure}[t]{0.48\textwidth}
    \centering
    \includegraphics[width=\textwidth]{Figs/reverse_probe_coin_k2_train.png}
    \caption{Train mode, $k = 2$ ($K_{\text{min}} = 4$)}
  \end{subfigure}
  \caption{Validation~$R^2$ of the reverse linear probe
  (predictor $\to$ hidden state) across layers~0--5, with $N = 2^{13}$
  sequences and positions $40$--$79$.
  Panels~(a)--(c): major-task sampling mode (posterior has $K_{\text{maj}}$
  components).
  Panel~(d): training-mixture mode (posterior has all $K_{\text{tot}}$
  components, data include minor tasks).
  The logits baseline is near zero in the major-only setting but rises sharply
  in the training mixture, where the model must encode task identity in its
  output distribution.}
  \label{fig:reverse_probe}
\end{figure}

\paragraph{Posterior-plane trajectory projection.}
To visualise how the model's internal representation of task identity evolves
over a sequence, we project hidden-state trajectories onto a two-dimensional
plane derived from the posterior--hidden linear fit.

\emph{Projection basis.}
From the reverse linear probe (fitted in major mode at late positions
$t \in \{100, \ldots, 127\}$), we obtain the weight matrix
$\mW \in \R^{K_{\mathrm{maj}} \times D}$ mapping the major-task posterior
to hidden space.
We define the final task vectors as the centred rows
$\bm{\tau}_j = \mW_j - \bar{\mW}$,
$j \in \{1, \ldots, K_{\mathrm{maj}}\}$,
where $\bar{\mW} = \frac{1}{K_{\mathrm{maj}}}\sum_j \mW_j$,
rescaled so that their Frobenius norm matches the actual hidden deviations at
the final time step.

\emph{Trajectory computation.}
For each task~$k$ (major, minor, or OOD) and time step~$t$, we compute the
batch-averaged hidden deviation
$\bar{\vh}_{k,t} - \bar{\vmu}_t$
(where
$\bar{\vmu}_t = \frac{1}{K_{\mathrm{maj}}}\sum_{j=1}^{K_{\mathrm{maj}}}
  \bar{\vh}_{j,t}$
is the mean over the major tasks) and project it onto the
plane spanned by $\{\bm{\tau}_j\}$ via least-squares regression with the
constraint $\sum_j \lambda_j = 1$.
The $R^2$ of this fit measures how much of the hidden trajectory lies in the
posterior plane.
The resulting plot shows each task's trajectory in the projected 2D~space, with
major tasks converging to their respective reference stars and OOD/minor~tasks
converging to positions within the major-task simplex.
Marker size encodes~$R^2$ and displayed time points are selected via
farthest-point sampling on the projected coordinates to avoid visual
clustering.

\emph{Experiments.}
We evaluate at layer~4 with $B = 64$ batch samples per task.
\Figref{fig:traj_proj} shows two settings.
In the first ($k = -1$, no minor tasks, panel~a), we generate $32$ OOD tasks
(fresh Dirichlet draws).
The three major-task trajectories (blue shades) start near the centre and
converge cleanly to their respective reference stars, with $R^2$ values
growing from~${\sim}0.66$ to~${\sim}0.96$ along the trajectory.
The OOD trajectories (red, dashed) fan out across the interior of the
major-task simplex: each OOD task converges toward a convex combination of the
major-task endpoints determined by its posterior under the major-task pool.
The high final~$R^2$ confirms that the posterior plane captures the dominant
structure of the hidden representation.

In the second setting ($k = 11$, $K_{\mathrm{min}} = 2048$, panel~b), we
include 16~OOD and 16~minor tasks.
The major-task trajectories still converge to their reference stars, but the
$R^2$ values are markedly lower (${\sim}0.03$--$0.14$), indicating that the
posterior plane explains only a small fraction of the hidden-state variance.
The OOD (red) and minor (green, dotted) trajectories cluster tightly in the
centre of the simplex rather than spreading toward individual corners.
This reflects the richer task structure: with 2048~minor tasks the model's
representation encodes a high-dimensional posterior that cannot be well
approximated by a 3-dimensional major-task plane.

\begin{figure}[ht]
  \centering
  \begin{subfigure}[t]{0.48\textwidth}
    \centering
    \includegraphics[width=\textwidth]{Figs/traj_proj_coin_k-1.png}
    \caption{$k = -1$ (no minor tasks), 32 OOD tasks}
  \end{subfigure}
  \hfill
  \begin{subfigure}[t]{0.48\textwidth}
    \centering
    \includegraphics[width=\textwidth]{Figs/traj_proj_coin_k11.png}
    \caption{$k = 11$ ($K_{\text{min}} = 2048$), 16 OOD + 16 minor}
  \end{subfigure}
  \caption{Posterior-plane trajectory projection at layer~4 ($B = 64$).
  Major-task trajectories (blue) converge to reference stars~($\bigstar$);
  OOD trajectories (red, dashed) and minor trajectories (green, dotted)
  settle within the simplex.
  Marker size encodes~$R^2$.
  With few tasks~(a) the posterior plane captures most variance ($R^2$ up
  to~$0.96$); with many minor tasks~(b) the plane explains little
  ($R^2 \sim 0.03$--$0.14$), reflecting a higher-dimensional posterior.}
  \label{fig:traj_proj}
\end{figure}

\paragraph{Comparing projected coefficients with the Bayesian posterior.}
To directly assess whether the model's hidden representation tracks the
Bayesian posterior, we compare the projection coefficients
$\bm{\lambda}(t) = (\lambda_1(t), \ldots, \lambda_{K_{\mathrm{maj}}}(t))$
from the trajectory projection with the true Bayesian posterior
$\P(Z = k \mid x_{1:t})$ computed from the known task distributions.

\emph{Computing $\bm{\lambda}(t)$.}
For a single sequence sample~$b$, let $\vd(t) = \vh_{k,t}^{(b)} - \bar{\vmu}_t
\in \R^D$ denote the hidden deviation at time~$t$, where $\bar{\vmu}_t$ is the
major-task mean as before.
We seek coefficients $\bm{\lambda}(t)$ satisfying
$\vd(t) \approx \sum_{j=1}^{K_{\mathrm{maj}}} \lambda_j(t)\, \bm{\tau}_j$
subject to $\sum_j \lambda_j = 1$.
With $K_{\mathrm{maj}} = 3$, the sum-to-one constraint is enforced by solving
the reduced system: let $\mX = [\bm{\tau}_1 \mid \bm{\tau}_2] \in \R^{D
\times 2}$ (dropping the last task vector since the basis is zero-mean).
We compute $\hat{\bm{w}} = (\mX^\top \mX)^{-1} \mX^\top \vd(t) \in \R^2$
via least-squares, set $\lambda_j = \hat{w}_j$ for $j \in \{1,2\}$, and
recover $\lambda_3 = 1 - \lambda_1 - \lambda_2$.
A uniform shift $s = (1 - \sum_j \hat{w}_j) / K_{\mathrm{maj}}$ is added to
all entries to ensure the constraint holds exactly.

Since the unconstrained $\bm{\lambda}$ may fall outside the probability
simplex, we project it onto $\Delta^{K_{\mathrm{maj}}-1}$ using the Euclidean
simplex projection of~\citet{duchi2008efficient}:
given $\bm{\lambda} \in \R^n$, sort in decreasing order, find
$\rho = \max\bigl\{j : \lambda_{(j)} - \tfrac{1}{j}
  \bigl(\sum_{i=1}^{j}\lambda_{(i)} - 1\bigr) > 0\bigr\}$,
set $\theta = \tfrac{1}{\rho}\bigl(\sum_{i=1}^{\rho}\lambda_{(i)} - 1\bigr)$,
and return $\max(\lambda_i - \theta, 0)$.

Since the reverse linear probe uses posterior probabilities (not
log-probabilities), the ideal relationship is
$\lambda_k(t) = \P(Z{=}k \mid x_{1:t})$ when the hidden state perfectly
encodes the posterior via the learned linear map.

\emph{Experiments.}
We evaluate at layer~5 on the $k = -1$ model (no minor tasks, $K_{\mathrm{tot}}
= 3$).
For each chosen task, 3~random sequence samples are drawn from the $B = 64$
batch.
For each sample, $\bm{\lambda}(t)$ is computed at every position and the true
Bayesian posterior $\P(Z{=}k \mid x_{1:t})$ is calculated from the raw token
sequence using the known major-task distributions.
The result is displayed as a $3 \times 3$ grid: one row per sample, one column
per task component~$k$.
Each cell overlays $\lambda_k(t)$ (solid) and $\P(Z{=}k \mid x_{1:t})$
(dashed) versus position~$t$.

\Figref{fig:lambda_vs_posterior} shows two settings.
When the generating task is a major task (task~1, panel~a), all three samples
converge rapidly: the projected coefficients closely track the true posterior,
with $\lambda_1(t)$ rising to~1 and $\lambda_2(t), \lambda_3(t)$ falling to~0
within the first ${\sim}20$ positions.
The solid and dashed curves are nearly indistinguishable after position~30,
confirming that the layer-5 hidden state faithfully encodes the Bayesian
posterior via the linear map~$\mW$.

When the generating task is an OOD task (task~5, panel~b), the picture changes
qualitatively.
Here the data are generated from a fresh Dirichlet draw $\vp_{\mathrm{ood}}$
that does not belong to the major-task pool.
The posterior~\eqref{eq:coin_posterior} is computed as if the true task were one
of the $K_{\mathrm{maj}} = 3$ major tasks, with a uniform prior over this pool:
\[
  \P(Z = k \mid x_{1:t})
  = \frac{\prod_{v} p_k(v)^{\,n_v^{(t)}}}
         {\sum_{j=1}^{3} \prod_{v} p_j(v)^{\,n_v^{(t)}}},
  \qquad k \in \{1, 2, 3\}.
\]
This is a \emph{misspecified} Bayesian inference: the likelihood model assumes
the data come from one of the major tasks, but the true generating distribution
$\vp_{\mathrm{ood}}$ lies outside this family.
Consequently the posterior need not concentrate on a single task;
it can oscillate as the observed token counts happen to favour different
major tasks at different positions.

The projected coefficients track this uncertainty:
$\bm{\lambda}(t)$ oscillates between major-task corners, closely following
the misspecified posterior dynamics.
The agreement between solid and dashed curves persists even in this
out-of-distribution regime, demonstrating that the model's hidden
representation maintains a faithful encoding of the (misspecified) Bayesian
posterior---not only for in-distribution tasks but also for novel tasks not
seen during training.

\begin{figure}[ht]
  \centering
  \begin{subfigure}[t]{\textwidth}
    \centering
    \includegraphics[width=\textwidth]{Figs/lambda_vs_posterior_task0.png}
    \caption{Major task (task~1): $\bm{\lambda}(t)$ converges quickly to the
    true posterior.}
  \end{subfigure}

  \vspace{0.5em}
  \begin{subfigure}[t]{\textwidth}
    \centering
    \includegraphics[width=\textwidth]{Figs/lambda_vs_posterior_task4.png}
    \caption{OOD task (task~5): $\bm{\lambda}(t)$ tracks the fluctuating
    posterior throughout the sequence.}
  \end{subfigure}
  \caption{Projected coefficients $\lambda_k(t)$ (solid) versus true Bayesian
  posterior $\P(Z{=}k \mid x_{1:t})$ (dashed) at layer~5, for the $k = -1$
  model (no minor tasks).
  Each row shows a different random sequence sample; each column shows one of
  the 3~major-task components.
  Close agreement confirms that the hidden representation faithfully encodes
  the posterior via the linear map~$\mW$.}
  \label{fig:lambda_vs_posterior}
\end{figure}

\paragraph{Quantifying the agreement between $\bm{\lambda}$ and the posterior.}
To summarise the agreement between the projected coefficients and the Bayesian
posterior across all tasks and samples, we compute three complementary metrics
at each position~$t$, each averaged over the $B$ batch samples and then over
tasks within each group (major or OOD).

\emph{(i) Total-variation distance.}
\begin{equation}
\label{eq:tv}
  \mathrm{TV}(t)
  = \frac{1}{2} \sum_{j=1}^{K_{\mathrm{maj}}}
    \bigl|\lambda_j(t) - p^*_j(t)\bigr|,
\end{equation}
where $p^*_j(t) = \P(Z{=}j \mid x_{1:t})$.
This measures the point-wise absolute discrepancy between the two
distributions; $\mathrm{TV} = 0$ when they coincide.

\emph{(ii) Cosine similarity.}
\begin{equation}
\label{eq:cos_sim}
  \cos\bigl(\bm{\lambda}(t),\, \vp^*(t)\bigr)
  = \frac{\bm{\lambda}(t) \cdot \vp^*(t)}
         {\|\bm{\lambda}(t)\| \;\|\vp^*(t)\|}.
\end{equation}
This captures directional agreement: two distributions that point toward the
same vertex of the simplex have cosine similarity near~1, regardless of small
absolute differences.

\emph{(iii) Rolling Pearson correlation.}
For each posterior component $j \in \{1, \ldots, K_{\mathrm{maj}}\}$, we
treat $\lambda_j(s)$ and $p^*_j(s)$ as paired time series and compute the
Pearson correlation within a sliding window of $w$~positions centred at~$t$:
\begin{equation}
\label{eq:rolling_corr}
  r_j(t) = \frac{
    \sum_{s \in W_t} \bigl(\lambda_j(s) - \bar{\lambda}_j\bigr)
                     \bigl(p^*_j(s)     - \bar{p}^*_j\bigr)
  }{
    \sqrt{\sum_{s \in W_t} \bigl(\lambda_j(s) - \bar{\lambda}_j\bigr)^2
    \;\cdot\;
    \sum_{s \in W_t} \bigl(p^*_j(s) - \bar{p}^*_j\bigr)^2}
  },
\end{equation}
where $W_t = \{t - w + 1, \ldots, t\}$, and $\bar{\lambda}_j$, $\bar{p}^*_j$
are the windowed means.
The final metric is $\bar{r}(t) = \frac{1}{K_{\mathrm{maj}}}
\sum_j r_j(t)$, averaged over components.
When both signals are effectively constant within the window
(standard deviation below a numerical threshold), that component is excluded
from the average, avoiding spurious noise.
We use $w = 20$ throughout.
High $\bar{r}(t)$ indicates that the two time series co-move within the
window---$\bm{\lambda}$ tracks the temporal dynamics of the posterior.

\paragraph{OOD and minor task metrics across training.}
To track how the model's internal task representations develop during training,
we evaluate two complementary projection bases at each training checkpoint and
each layer.

\emph{Two projection bases.}
The \emph{major basis} consists of the centred rows of $\mW$ from the reverse
linear probe (re-fitted at each checkpoint, since the model's representations
evolve during training), defining the $K_{\mathrm{maj}}$-dimensional posterior
subspace.
The \emph{minor basis} consists of the top-$K_{\mathrm{maj}}$ right singular
vectors of the minor tasks' batch-averaged hidden deviations at the final
time step.
For each basis, we compute the mean $R^2$ at the final time step for OOD~tasks
($R^2_{\mathrm{OOD}}$: how well OOD hidden states lie in the projection plane)
and for minor tasks ($R^2_{\mathrm{minor}}$).

\emph{Experiment details.}
We evaluate at layer~3 on 6~experiments with
$k \in \{0, 4, 6, 8, 10, 12\}$ (i.e.\ $K_{\mathrm{min}} \in \{1, 16, 64,
256, 1024, 4096\}$), sweeping training steps from~0 to~30\,000 at intervals
of~100 ($B = 64$, $N_{\mathrm{ood}} = 30$, $N_{\mathrm{minor}} = 64$).
An OOD-loss curve from each experiment's training log is shown alongside.

\Figref{fig:ood_minor_training} shows the four $R^2$ metrics and the OOD loss.

\emph{Major basis, OOD tasks} (panel~a, \texttt{maj\_r2\_ood}).
When $k = 0$ (essentially no minor tasks), the major-basis $R^2$ for OOD tasks
remains high (${\sim}0.6$--$0.7$) throughout training: the 3-dimensional
posterior plane captures most of the OOD hidden-state variance.
As the minor-task pool grows, the $R^2$ drops steadily (${\sim}0.1$ for
$k = 12$), indicating that the model allocates representational capacity to
encode the richer task structure, and the low-dimensional major-task plane
becomes insufficient.

\emph{Major basis, minor tasks} (panel~b, \texttt{maj\_r2\_min}).
A similar pattern holds: the major-basis $R^2$ for minor tasks starts high and
declines during training for all~$k$.
For $k = 4$ it drops to nearly~0, indicating that with 16~minor tasks the
model learns representations that are poorly explained by the 3-task posterior
plane.

\emph{Minor basis, minor tasks} (panel~c, \texttt{min\_r2\_min}).
Here the top-3 SVD basis of the minor tasks' own hidden deviations is used.
A striking pattern emerges: experiments with \emph{more} minor tasks achieve
\emph{higher} $R^2$ (${\sim}0.78$ for $k = 12$ vs.\ ${\sim}0.63$ for
$k = 4$).
This suggests that as the number of minor tasks grows, the model's hidden
representations for minor tasks increasingly collapse into a low-dimensional
subspace---a top-3 SVD basis suffices to capture a larger share of the
variance.
Intuitively, when $K_{\mathrm{min}}$ is large, the model cannot allocate a
separate direction for each minor task;
instead it compresses the minor-task variation into a compact subspace whose
dimensionality grows much more slowly than the number of tasks.

\emph{Minor basis, OOD tasks} (panel~d, \texttt{min\_r2\_ood}).
The same trend appears for OOD tasks projected onto the minor basis:
larger $K_{\mathrm{min}}$ yields higher $R^2$, confirming that the
minor-task subspace also captures OOD variation.
For $k = 12$, the $R^2$ stabilises around~$0.80$, whereas for $k = 4$ it
settles near~$0.55$.

\emph{OOD loss} (panel~e).
The OOD next-token prediction loss decreases monotonically with~$k$: models
trained with more minor tasks achieve better OOD generalisation
(${\sim}1.81$ for $k = 12$ vs.\ ${\sim}1.97$ for $k = 4$), confirming that
exposure to a richer task pool during training improves out-of-distribution
performance.

\emph{Remark on the minor basis.}
The increasing $R^2$ of the top-3 SVD basis with~$K_{\mathrm{min}}$ suggests
that a fixed low-rank subspace is a reasonable summary of minor-task
representations when the task pool is large.
However, a rank-3 basis is chosen here only for comparability with the
$K_{\mathrm{maj}} = 3$ major basis, and there is no reason to expect it to be
optimal.
A natural next step is to study how the effective dimensionality of the
minor-task subspace scales with~$K_{\mathrm{min}}$---for instance by examining
the singular-value spectrum or the stable rank of the minor tasks' hidden
deviation matrix---and to determine whether the compression into a
low-dimensional subspace is a consequence of the model's architecture, the
task structure, or both.

\begin{figure}[ht]
  \centering
  \begin{subfigure}[t]{0.48\textwidth}
    \centering
    \includegraphics[width=\textwidth]{Figs/training_maj_r2_ood.png}
    \caption{\texttt{maj\_r2\_ood}: major basis, OOD tasks}
  \end{subfigure}
  \hfill
  \begin{subfigure}[t]{0.48\textwidth}
    \centering
    \includegraphics[width=\textwidth]{Figs/training_maj_r2_min.png}
    \caption{\texttt{maj\_r2\_min}: major basis, minor tasks}
  \end{subfigure}

  \vspace{0.5em}
  \begin{subfigure}[t]{0.48\textwidth}
    \centering
    \includegraphics[width=\textwidth]{Figs/training_min_r2_min.png}
    \caption{\texttt{min\_r2\_min}: minor basis, minor tasks}
  \end{subfigure}
  \hfill
  \begin{subfigure}[t]{0.48\textwidth}
    \centering
    \includegraphics[width=\textwidth]{Figs/training_min_r2_ood.png}
    \caption{\texttt{min\_r2\_ood}: minor basis, OOD tasks}
  \end{subfigure}

  \vspace{0.5em}
  \begin{subfigure}[t]{0.48\textwidth}
    \centering
    \includegraphics[width=\textwidth]{Figs/training_ood_loss.png}
    \caption{OOD next-token prediction loss}
  \end{subfigure}
  \caption{OOD and minor task metrics across training at layer~3, for
  $k \in \{0, 4, 6, 8, 10, 12\}$ (one curve per experiment).
  Panels~(a)--(d) show the $R^2$ of projecting OOD or minor task
  representations onto either the major (W-based) or minor (SVD-based) basis.
  Panel~(e) shows the OOD loss.
  As $K_{\text{min}}$ grows, the major-basis $R^2$ decreases (the posterior
  plane becomes insufficient) while the minor-basis $R^2$ \emph{increases}
  (minor-task representations collapse into a low-dimensional subspace).}
  \label{fig:ood_minor_training}
\end{figure}
